% ============================================================
\setpartpreamble{\partintro
  {Which building blocks belong in a shared core, which in provider packages, and how a second
  provider is built on the same core.}
  {
    \item Chapter~\ref{ch:what-belongs-in-the-core}: What belongs in the core
    \item Chapter~\ref{ch:a-second-provider-sql-server}: A second provider: SQL Server
  }}
\part{Multiple providers}\label{part:providers}

% ------------------------------------------------------------
\chapter{What belongs in the core}\label{ch:what-belongs-in-the-core}
\versionrisk{stable}

A second provider is the acid test for the package architecture. It shows which parts are
truly provider-independent, and the first impression is misleading: a building block can know
no SQL syntax and no provider type and still need a relational database, because it works
with ADO.NET connections or with tables. The sample therefore has three layers, each in a
package of its own:

\begin{datatable}{H{1.21}H{0.45}H{1.34}}
  \textbf{Building block} & \textbf{Layer} & \textbf{Reason} \\
  \tablehead
  \code{Tenant\allowbreak{}Id}, converter, JSON reader/writer & Core & Pure .NET types and
  EF value converters \\
  \tablerule
  \code{ITenant\allowbreak{}Owned}, \code{ITenant\allowbreak{}Provider}, attributes & Core &
  For entity assemblies without a provider dependency \\
  \tablerule
  \code{Is\allowbreak{}Tenant\allowbreak{}Property(\allowbreak{})}, candidate lookup & Core
  & Metadata only \\
  \tablerule
  Runtime annotation, \code{Get\allowbreak{}Tenant\allowbreak{}Property\allowbreak{}Name} &
  Core & Read by the provider-neutral interceptors \\
  \tablerule
  SaveChanges interceptor & Core & Works only with the change tracker \\
  \tablerule
  Protection of bulk updates (query expression interceptor) & Core & Only inspects the LINQ
  expression tree \\
  \tablerule
  Evaluatable filter with attribute & Core & Recognizes function classes of all packages \\
  \tablerule
  Check while loading (materialization interceptor) & Core & Only reads the model, the
  context and the row's values \\
  \tablerule
  Value generation, logging, \code{Add\allowbreak{}Acme\allowbreak{}Tenancy} & Core & No
  database involved \\
  \tablerule
  \code{Tenant\allowbreak{}Session\allowbreak{}Interceptor} (abstract) & Relational & Works
  on the \code{Db\allowbreak{}Connection}; only the SQL command differs per provider \\
  \tablerule
  \code{Tenant\allowbreak{}Table\allowbreak{}Mapping\allowbreak{}Validator} & Relational &
  Tables and columns: TPT, TPC, entity splitting \\
  \tablerule
  \code{Database.\allowbreak{}Apply\allowbreak{}Tenant\allowbreak{}Session(\allowbreak{})} &
  Relational & Uses the session interceptor of the provider \\
  \tablerule
  Type mapping, translators, migrations, RLS conventions & Provider & Store types, SQL
  syntax, schema objects \\
  \tablerule
  Custom SQL nodes, aggregates & Provider & Depend on syntax and provider factory \\
\end{datatable}

The rule of thumb therefore has three stages: \textbf{whatever needs no database belongs in
the core; whatever needs a relational database but no SQL dialect belongs in the relational
layer; SQL syntax and store types belong in the provider.} The core package references only
\code{Microsoft.\allowbreak{}Entity\allowbreak{}Framework\allowbreak{}Core}, the package
\code{Acme.\allowbreak{}Entity\allowbreak{}Framework\allowbreak{}Core.\allowbreak{}Tenancy.\allowbreak{}Relational}
adds \code{Microsoft.\allowbreak{}Entity\allowbreak{}Framework\allowbreak{}Core.\allowbreak{}Relational},
and both provider packages build on it. \code{Package\allowbreak{}Layer\allowbreak{}Tests}
makes sure that no relational type creeps back into the core.

\begin{note}[Separate the relational layer from the start]
  If you might support a non-relational provider such as Cosmos DB later, put the relational
  layer into a package of its own from the beginning. Moving types into another assembly
  afterwards breaks binary compatibility, and entity assemblies that only referenced the core
  lose the transitive dependency on
  \code{Microsoft.\allowbreak{}Entity\allowbreak{}Framework\allowbreak{}Core.\allowbreak{}Relational}
  they may have relied on for \code{To\allowbreak{}Table} or
  \code{Has\allowbreak{}Column\allowbreak{}Type}. The sample made exactly this move in
  edition~1.2.
\end{note}

\begin{important}[Name collisions between provider packages]
  Two provider packages that both define an extension method
  \code{Use\allowbreak{}Acme\allowbreak{}Tenancy} on
  \code{Db\allowbreak{}Context\allowbreak{}Options\allowbreak{}Builder} cannot be referenced
  at the same time: the call would be ambiguous. The PostgreSQL package uses EF's builder,
  the SQL Server package the provider's builder
  (\code{Use\allowbreak{}Sql\allowbreak{}Server(\allowbreak{}cs,\allowbreak{} x =>
  x.\allowbreak{}Use\allowbreak{}Acme\allowbreak{}Tenancy(\allowbreak{}))}). If you plan
  several provider packages, use the provider builder consistently from the start, as the
  NetTopologySuite plugin does.
\end{important}

% ------------------------------------------------------------
\chapter{A second provider: SQL Server}\label{ch:a-second-provider-sql-server}
\versionrisk{medium}

The package
\code{Acme.\allowbreak{}Entity\allowbreak{}Framework\allowbreak{}Core.\allowbreak{}Tenancy.\allowbreak{}Sql\allowbreak{}Server}
implements the same idea with SQL Server's means. The differences show what a plugin really
has to solve anew per database:

\begin{datatable}{H{0.81}H{1.00}H{1.19}}
  \textbf{Aspect} & \textbf{PostgreSQL} & \textbf{SQL Server} \\
  \tablehead
  Tenant in the session & \code{set\_\allowbreak{}config(\allowbreak{}name,\allowbreak{}
  value,\allowbreak{} false)} &
  \code{sp\_\allowbreak{}set\_\allowbreak{}session\_\allowbreak{}context} \\
  \tablerule
  \code{Current\allowbreak{}Tenant(\allowbreak{})} in SQL &
  \code{nullif(\allowbreak{}current\_\allowbreak{}setting(\allowbreak{}.\allowbreak{}.\allowbreak{}.\allowbreak{}),\allowbreak{}
  '')::uuid} & \code{CAST(\allowbreak{}SESSION\_\allowbreak{}CONTEXT(\allowbreak{}N'…') AS
  uniqueidentifier)} \\
  \tablerule
  Column type & \code{uuid} & \code{uniqueidentifier} \\
  \tablerule
  RLS mechanism & \code{CREATE POLICY} per table & Predicate function + \code{CREATE
  SECURITY POLICY} \\
  \tablerule
  Read / write & \code{USING} / \code{WITH CHECK} & \code{FILTER} / \code{BLOCK
  .\allowbreak{}.\allowbreak{}.\allowbreak{} AFTER INSERT/\allowbreak{}UPDATE} \\
  \tablerule
  Table owner & Only bound with \code{FORCE} & No FORCE; predicates are part of the query \\
  \tablerule
  \code{DROP TABLE}, renaming & Policies disappear with it or follow & Policy must go first and be recreated afterwards (SCHEMABINDING) \\
  \tablerule
  DDL in batches & Transactional, freely combinable & \code{CREATE
  FUNCTION/\allowbreak{}SCHEMA} alone in a batch:
  \code{EXEC(\allowbreak{}N'.\allowbreak{}.\allowbreak{}.\allowbreak{}')} \\
  \tablerule
  Table mappings & \multicolumn{2}{>{\RaggedRight\arraybackslash}p{0.6\linewidth}}{Same
  contract from the relational package (chapter~\ref{ch:conventions}, page~\pageref{sec:supported-table-mappings})} \\
  \tablerule
  Deliberate exception & Attribute or
  \code{Has\allowbreak{}No\allowbreak{}Row\allowbreak{}Level\allowbreak{}Security(\allowbreak{})}
  & Attribute only \\
\end{datatable}

\section{Activation through the provider builder}

\begin{sourcecode}[Usage and Infrastructure/\allowbreak{}Sql\allowbreak{}Server\allowbreak{}Tenancy\allowbreak{}Options\allowbreak{}Builder\allowbreak{}Extensions.\allowbreak{}cs (\allowbreak{}shortened, simplified)]{csharp}
services.AddAcmeTenancy<HttpTenantProvider>();                    // core package
services.AddDbContext<ShopContext>((sp, o) => o
    .UseSqlServer(connectionString, sql => sql.UseAcmeTenancy(sp)));

public static SqlServerDbContextOptionsBuilder UseAcmeTenancy(
    this SqlServerDbContextOptionsBuilder sqlServerOptionsBuilder,
    string sessionKey = "TenantId",
    string policySchema = "acme")
{
    // From the provider builder to the actual DbContextOptionsBuilder
    var optionsBuilder =
        ((IRelationalDbContextOptionsBuilderInfrastructure)sqlServerOptionsBuilder)
        .OptionsBuilder;

    var extension = optionsBuilder.Options.FindExtension<SqlServerTenancyOptionsExtension>()
                    ?? new SqlServerTenancyOptionsExtension();
    // ... validate and set settings ...
    ((IDbContextOptionsBuilderInfrastructure)optionsBuilder)
        .AddOrUpdateExtension(extension.WithSettings(settings));

    return sqlServerOptionsBuilder;
}
\end{sourcecode}

\section{Migrations: security policy instead of policy expression}

For SQL Server, the convention only records the tenant column; the SQL is produced in the
generator. The predicate function is the same for all tables and is created idempotently:

\begin{sourcecode}[Migrations/\allowbreak{}Internal/\allowbreak{}Sql\allowbreak{}Server\allowbreak{}Tenancy\allowbreak{}Migrations\allowbreak{}Sql\allowbreak{}Generator.\allowbreak{}cs (\allowbreak{}excerpt, simplified)]{csharp}
protected virtual void CreatePolicy(string table, string? schema, string column,
    MigrationCommandListBuilder builder)
{
    var function = Delimit(PredicateFunctionName, Settings.PolicySchema);
    var target = Delimit(table, schema);
    var predicate = $"{function}({Delimit(column)})";

    // CREATE SCHEMA/FUNCTION must be alone in a batch -> EXEC(N'...')
    Statement(builder, $"IF SCHEMA_ID({Literal(Settings.PolicySchema)}) IS NULL "
        + $"EXEC({Literal($"CREATE SCHEMA {Delimit(Settings.PolicySchema)}")})");

    var functionSql =
        $"CREATE FUNCTION {function}(@TenantId uniqueidentifier) "
        + "RETURNS TABLE WITH SCHEMABINDING AS RETURN SELECT 1 AS [Allowed] "
        + $"WHERE @TenantId = CAST(SESSION_CONTEXT({Literal(Settings.SessionKey)}) "
        + "AS uniqueidentifier)";
    Statement(builder, $"IF OBJECT_ID({Literal(function)}, N'IF') IS NULL "
        + $"EXEC({Literal(functionSql)})");

    Statement(builder,
        $"CREATE SECURITY POLICY {PolicyName(table, schema ?? "dbo")} "
        + $"ADD FILTER PREDICATE {predicate} ON {target}, "
        + $"ADD BLOCK PREDICATE {predicate} ON {target} AFTER INSERT, "
        + $"ADD BLOCK PREDICATE {predicate} ON {target} AFTER UPDATE "
        + "WITH (STATE = ON)");
}

protected override void Generate(DropTableOperation operation, IModel? model,
    MigrationCommandListBuilder builder, bool terminate = true)
{
    // SCHEMABINDING: the policy first, found by the table it binds
    DropPolicy(operation.Name, TableSchema(operation.Schema, model), builder);
    base.Generate(operation, model, builder, terminate);
}
\end{sourcecode}

\code{SCHEMABINDING} only accepts two-part names. A table without an explicit schema is
therefore qualified by the callers with the model's default schema or, failing that, with
\code{dbo}; otherwise SQL Server rejects the security policy.

\begin{important}[Nested SQL]
  The \code{EXEC(\allowbreak{}N'.\allowbreak{}.\allowbreak{}.\allowbreak{}')} contains SQL
  in a string literal that in turn contains literals. The nesting is produced exclusively
  through the string mapping's \code{Generate\allowbreak{}Sql\allowbreak{}Literal}, which
  doubles quotes correctly; identifiers go through \code{Delimit\allowbreak{}Identifier}.
  Session key and schema are additionally validated during configuration.
\end{important}

\section{Schema changes on protected tables}\label{sec:sqlserver-schema-bound-changes}

\code{SCHEMABINDING} protects the predicate function, but it also binds the table. SQL Server
rejects \code{sp\_\allowbreak{}rename} for the table and the tenant column, \code{ALTER SCHEMA
.\allowbreak{}.\allowbreak{}.\allowbreak{} TRANSFER}, \code{ALTER COLUMN} and \code{DROP COLUMN}
as long as a security policy binds the table: “object participates in enforced dependencies”.
This cannot be solved per operation. EF orders a migration's operations itself, and when the
tenant column is renamed, the \code{Rename\allowbreak{}Column\allowbreak{}Operation} comes
\emph{before} the \code{Alter\allowbreak{}Table\allowbreak{}Operation} that carries the changed
annotation. A plain table rename produces no \code{Alter\allowbreak{}Table\allowbreak{}Operation}
at all.

The generator therefore overrides the entry point that receives the whole operation list. It
finds all tables whose bound objects change, follows their renames and wraps the list. At the
front it removes the policies of the affected tables under their old names. At the end it
recreates them with the final table, schema and column names from the target model:

\begin{sourcecode}[Generated SQL when renaming the tenant column (\allowbreak{}shortened)]{sql}
DECLARE @policy nvarchar(400) = (SELECT TOP (1) ... FROM sys.security_policies AS p
    WHERE EXISTS (SELECT 1 FROM sys.security_predicates AS s WHERE s.object_id = p.object_id
        AND s.target_object_id = OBJECT_ID(N'[dbo].[Invoice]')
        AND s.predicate_definition LIKE N'%[[]fn[_]tenant[_]predicate]%'));
IF @policy IS NOT NULL EXEC(N'DROP SECURITY POLICY ' + @policy);
EXEC sp_rename N'[Invoice].[TenantId]', N'Mandant', 'COLUMN';
EXEC sp_rename N'[Invoice].[IX_Invoice_TenantId]', N'IX_Invoice_Mandant', 'INDEX';
IF SCHEMA_ID(N'acme') IS NULL EXEC(N'CREATE SCHEMA [acme]');
IF OBJECT_ID(N'[acme].[fn_tenant_predicate]', N'IF') IS NULL EXEC(N'CREATE FUNCTION ...');
CREATE SECURITY POLICY [acme].[tenant_dbo_Invoice]
    ADD FILTER PREDICATE [acme].[fn_tenant_predicate]([Mandant]) ON [dbo].[Invoice], ...
\end{sourcecode}

Three decisions are built into this:

\begin{itemize}
  \item \textbf{A transaction instead of switching off.} EF runs a migration's commands in a
    transaction, and DDL is transactional in SQL Server. If a step fails, SQL Server also
    rolls back the removal of the policy. A permanently removed policy, or one switched off
    with \code{STATE = OFF}, cannot come about this way. The Down path is the same wrapping
    with the names swapped.
  \item \textbf{Removal by binding, not by name.} Only a policy whose predicate uses the
    plugin's predicate function and that binds exactly this table is dropped. Policies of
    others and tables without RLS stay untouched, and policies under an earlier name are
    found all the same.
  \item \textbf{Schema-qualified names.} Up to 1.0 the policy was called
    \code{tenant\_\allowbreak{}Table}, so tables with the same name in two schemas collided.
    New policies are called \code{tenant\_\allowbreak{}Schema\_\allowbreak{}Table}. Existing
    installations keep the old name until the table is touched the next time; then the
    generator replaces the policy. The remaining ambiguity (schema \code{a\_b} with table
    \code{c} versus schema \code{a} with table \code{b\_c}) and names longer than 128
    characters are rejected by the convention while the model is built.
\end{itemize}

Limits: restoring needs the target model. That is the normal case, since
\code{dotnet ef migrations add} creates a designer file for every migration. If it is
missing, or a changed table does not appear in it, the generator stops instead of dropping a
policy. The application's own schema-bound objects (views with \code{SCHEMABINDING}, its own
security policies) are not handled by the plugin; they still block the same operations. Data
operations that EF sorts between the schema changes run inside the transaction without a
policy, that is with the rights of the migration role.
\code{Sql\allowbreak{}Server\allowbreak{}Migration\allowbreak{}Tests} checks with diffs
generated by EF against a real SQL Server: create, rename column and table, move schema, every
step back and drop, with read and write isolation through the app login after every step. In
addition: tables with the same name in two schemas, a policy under the old name and someone
else's policy on a table without RLS.

\section{Changing the session key and the policy schema}\label{sec:sqlserver-session-settings}

\code{session\allowbreak{}Key} and \code{policy\allowbreak{}Schema} are options, but they act
in two places: at run time in the session interceptor, which sets the key, and in the database in
the predicate function, which reads it. Up to 1.0 the generator read both values from the
\emph{current} options and created the function only if it was missing. Switching from
\code{TenantId} to \code{NewTenant} therefore produced no migration. The interceptor then set the
new key, the function kept reading the old one, and no tenant saw any rows. Worse still: run
again, for example for a new test database, an old migration would have created objects with the
new values.

Since 1.1 the package separates the two roles. The convention writes key and schema into the
model as the model annotations
\code{Acme:\allowbreak{}Sql\allowbreak{}Server\allowbreak{}Rls:\allowbreak{}Session\allowbreak{}Key}
and \code{Acme:\allowbreak{}Sql\allowbreak{}Server\allowbreak{}Rls:\allowbreak{}Policy\allowbreak{}Schema},
and so into the snapshot. The annotation provider passes them on to the relational model, and the
differ produces an \code{Alter\allowbreak{}Database\allowbreak{}Operation} with the old and the
new value when they change. The generator reads both values from the target model of the
migration it is generating; the options only apply to models that do not record the values. That
way every migration stays with the state it was generated from.

\begin{datatable}{H{0.70}H{1.30}}
  \textbf{Case} & \textbf{Behavior of the migration} \\
  \tablehead
  Unchanged & No operation, no phantom change \\
  \tablerule
  Key or schema changed, old value in the snapshot & In one transaction: remove the policies of
  all protected tables, drop the old function, create the function with the new values,
  recreate the policies. The Down path does the same in reverse. \\
  \tablerule
  First migration of a context, or the first after the upgrade from 1.0 (old value unknown) &
  Nothing is dropped. An existing function must read the key that the model names, otherwise the
  migration stops with \code{THROW 50000}. \\
  \tablerule
  Down to a snapshot before 1.1 & Function and policies stay as they are. \\
\end{datatable}

The third row is a deliberate limit. Without a recorded old value there is no proof that an
existing function belongs to the plugin and not to another context. An existing installation
therefore changes its key in two steps: first the first migration after the upgrade with the key
unchanged, then the change in a migration of its own.

The same check protects several contexts in one database. There is one function per policy
schema, and every policy of the plugin checks the function's definition after creating it. If a
second context configures another key in the same schema, its migration fails instead of binding
its tables to a function that reads the wrong key. Contexts with different keys need different
policy schemas.

Run time and database can still drift apart briefly. If a new app version with a new key runs
before the migration has been applied, the interceptor sets a key that the function does not
read. The result is not a data leak but \emph{fail closed}: no rows, and writes fail on the BLOCK
predicate. The migration therefore belongs before the new app version.
\code{Sql\allowbreak{}Server\allowbreak{}Session\allowbreak{}Settings\allowbreak{}Tests} checks
without a database the annotations, the snapshot, the empty diff, the detected change and the
historical migration. Against a real SQL Server it checks changing key and schema with Up and
Down, the upgrade of a 1.0 installation and the second context.

\section{Run time: same flow, different command}

% source: src/Acme.EntityFrameworkCore.Tenancy.SqlServer/Runtime/SqlServerTenantSessionInterceptor.cs
\begin{sourcecode}[Runtime/\allowbreak{}Sql\allowbreak{}Server\allowbreak{}Tenant\allowbreak{}Session\allowbreak{}Interceptor.\allowbreak{}cs]{csharp}
public class SqlServerTenantSessionInterceptor : TenantSessionInterceptor
{
    public SqlServerTenantSessionInterceptor(ITenantProvider tenants)
        : base(tenants)
    {
    }

    protected override DbCommand CreateSessionCommand(
        DbConnection connection, DbContext? context, TenantId? tenant)
    {
        var key = context?.GetService<ISqlServerTenancySingletonOptions>()
                      .Settings.SessionKey
                  ?? SqlServerTenancySettings.Default.SessionKey;

        var command = connection.CreateCommand();
        command.CommandText = "EXEC sp_set_session_context @key = @key, @value = @value";
        AddParameter(command, "@key", key);
        AddParameter(command, "@value", tenant?.Value);    // NULL: no rows visible
        return command;
    }
}
\end{sourcecode}

\section{Deliberately not ported}

The SQL Server package is a self-contained but leaner provider. Fluent API policies, the optional
query filter, the snapshot code generator and scaffolding are missing there, and the custom
SQL node and the aggregate of the functions sample package exist for PostgreSQL only. That is a product decision, not a technical obstacle: each of these chapters
can be transferred to SQL Server with the same patterns. The list belongs in the package's
documentation so users know what to expect.

% source: tests/Acme.EntityFrameworkCore.Tenancy.SqlServer.Tests/SqlServerTests.cs
\begin{sourcecode}[tests/\allowbreak{}.\allowbreak{}.\allowbreak{}.\allowbreak{}Sql\allowbreak{}Server.\allowbreak{}Tests/\allowbreak{}Sql\allowbreak{}Server\allowbreak{}Tests.\allowbreak{}cs (\allowbreak{}excerpt)]{csharp}
[Fact]
public async Task Block_predicate_rejects_rows_of_other_tenant()
{
    await using var db = ContextFor(TenantId.New());

    var ex = await Assert.ThrowsAsync<SqlException>(() => db.Database.ExecuteSqlAsync(
        $"INSERT INTO [Orders] ([TenantId], [Total]) VALUES ({Guid.NewGuid()}, 1)"));

    Assert.Equal(33504, ex.Number);      // violation of a BLOCK PREDICATE
}
\end{sourcecode}

As with PostgreSQL, the tests run on two levels: without a database for model, query and
migration SQL, and with Testcontainers against a real SQL Server. There, the application
connects with its own login instead of \code{sa} to reproduce the production situation.
